% =====================================================================
%  Dissolution of the collinear structure in the solar-sail CR3BP
%  Venue-neutral preprint source.  To retarget:
%    JGCD / AIAA   -> \documentclass[]{aiaa-tc} or the AIAA template
%    CMDA / Acta   -> \documentclass{svjour3} / elsarticle
%  Only the preamble changes; the body is portable.
% =====================================================================
\documentclass[11pt,a4paper]{article}

\usepackage[margin=1in]{geometry}
\usepackage{amsmath,amssymb,bm}
\usepackage{graphicx}
\usepackage{booktabs}
\usepackage{siunitx}
\usepackage[hidelinks]{hyperref}
\usepackage[font=small,labelfont=bf]{caption}
\usepackage{authblk}

\sisetup{group-minimum-digits=4, detect-all}
\graphicspath{{../}}

\newcommand{\rhat}{\hat{\bm{r}}}
\newcommand{\phat}{\hat{\bm{p}}}
\newcommand{\qhat}{\hat{\bm{q}}}
\newcommand{\nhat}{\hat{\bm{n}}}
\newcommand{\nd}{\text{nd}}

\title{\bfseries Dissolution of the collinear structure in the
solar-sail circular restricted three-body problem:\\
an exact tidal-parity threshold and a universal frequency bound}

\author[1]{Bishwaswarup}
\affil[1]{\small Indian Institute of Science, Bangalore 560012, India\\
\texttt{bishwaswarup@iisc.ac.in}}

\date{\today}

\begin{document}
\maketitle

\begin{abstract}
\noindent
A face-on solar sail does not displace the collinear Lagrange point of the
circular restricted three-body problem; it dissolves it. Because the radiation
force is radial and therefore conservative, it rescales solar gravity as
$(1-\mu)\to(1-\beta)(1-\mu)$, and the on-axis equilibrium degenerates smoothly
into a heliocentric hovering point as the lightness number $\beta$ grows. We
characterise that degeneration and obtain two exact results. First, imposing
\emph{tidal parity} --- the condition that the smaller primary's contribution to
the linearised dynamics equal the centrifugal background, $s\equiv\mu/r_2^3=1$,
equivalently $r_2=\mu^{1/3}=3^{1/3}r_H$ --- on the on-axis force balance yields
the closed form
$\beta_{\mathrm{crit}}=1-\bigl(1-\mu^{1/3}\bigr)^2=2\mu^{1/3}-\mu^{2/3}$,
exact for every mass ratio and verified against independent root-finding to
\num{2.1e-16} over $\mu\in[10^{-7},10^{-2}]$. For Sun--Earth this gives
$\beta_{\mathrm{crit}}=0.0286$. Second, writing $c_2$ for the second Legendre
coefficient of the collinear expansion --- the quantity a sail detunes
continuously, and which is otherwise fixed by the mass ratio and the choice of
point --- the ratio of the out-of-plane to in-plane linear frequencies at any
collinear equilibrium is confined to $\nu/\omega\in[2\sqrt2/3,\,1)$, a band
\SI{5.72}{\percent} wide containing no low-order rational. The bound is attained
exactly at $c_2=8/5$, where the three linear rates stand in the ratio
$\lambda_u^2:\nu^2:\omega^2=7:8:9$. Both $c_2=8/5$ and the bound are independent
of $\mu$. That the two frequencies are nearly equal at collinear points has been
remarked upon~\cite{ceccaroni2016}; we make the statement exact. We emphasise that neither result is a bifurcation: the saddle persists
for all $\beta$ and the degeneration is smooth, so these are exact structural
landmarks rather than dynamical transitions, and we show by explicit search that
no commensurability exists anywhere to make one. Reducing $\beta$ from primary
mission specifications, we find that no solar sail yet flown exceeds
$\beta=0.0061$, placing tidal parity a factor of $4.7$ beyond demonstrated
capability. Finally, linearising the sail acceleration in its two attitude
angles shows that the face-on nominal is a \emph{singular} configuration for
station-keeping: $\partial\bm{a}/\partial\delta$ vanishes identically at
$\alpha_0=0$, the Kalman rank drops to $4/6$, and any non-zero cone angle
restores controllability.
\end{abstract}

\noindent\textbf{Keywords:} solar sail; restricted three-body problem;
artificial equilibria; collinear points; Legendre coefficient $c_2$;
Hill radius; halo orbits; controllability

% =====================================================================
\section{Introduction}
% =====================================================================

Artificial equilibria generated by continuous low thrust have been studied since
McInnes \emph{et al.}~\cite{mcinnes1994} mapped the surfaces of sail-displaced
equilibria in the restricted three-body problem, and the periodic orbits about
them are equally well established: Baoyin and McInnes~\cite{baoyin2006}
constructed halo orbits at the Sun--Earth artificial $L_1$, Waters and
McInnes~\cite{waters2007} treated families above the ecliptic, and Farr\'es and
Jorba~\cite{farres2010,farres2016} analysed the quasi-periodic dynamics and
station-keeping in the same region. The standard reference
remains~\cite{mcinnes1999}.

What that literature supplies is numerical: the equilibrium abscissa is obtained
by root-finding on the on-axis balance for each $\beta$, and the orbit families
by Lindstedt--Poincar\'e expansion or numerical continuation. The equilibrium
coordinates enter as computed inputs. What it does not supply is an algebraic
characterisation of \emph{how the collinear structure itself degrades} as $\beta$
grows --- at what sail performance the smaller primary ceases to organise the
dynamics at all, and whether that point admits a closed form.

This paper supplies that characterisation. Our starting observation is that a
face-on sail ($\alpha=0$) exerts a purely radial and hence conservative force,
so it does not perturb the collinear point: it rescales one primary's mass and
the equilibrium migrates toward a \emph{heliocentric hovering point} at which
reduced solar gravity balances centrifugal acceleration at the synchronous rate.
The three-body character of the equilibrium therefore does not shift, it decays,
and the useful question is the rate and the scale of that decay.

Section~\ref{sec:model} sets out the model and the conservative rescaling.
Section~\ref{sec:dissolution} defines the saddle strength $s=\mu/r_2^3$ and
measures its four-decade collapse. Section~\ref{sec:closed} derives the exact
tidal-parity threshold, the paper's principal result, and validates it across
five decades of mass ratio. Section~\ref{sec:freq} establishes the frequency
bound $\nu/\omega\in[2\sqrt2/3,1)$, locates its exact extremum at $c_2=8/5$, and
reports a systematic negative: no low-order resonance is attainable anywhere,
which is what forbids interpreting either threshold as a bifurcation.
Section~\ref{sec:families} tracks the halo families across the flown sail band
by pseudo-arclength continuation. Section~\ref{sec:control} derives the sail
control Jacobian and the singularity of the face-on nominal.
Section~\ref{sec:tech} places $\beta_{\mathrm{crit}}$ against every sail flown
to date. Section~\ref{sec:limits} states the limitations plainly.

\paragraph{On terminology.}
We use \emph{tidal parity} for the condition $s=1$. The term has recently
appeared in translunar astrodynamics for the analogous balance between
lunisolar perturbations and Earth gravity; we adopt it in that spirit and claim
no coinage. Likewise, the scale $\mu^{1/3}$ is not new --- it is the Hill radius
up to the factor $3^{1/3}$, and is standard in tidal-stability analysis. Our
contribution is not the scale but the exact sail parameter required to place a
spacecraft on it.

% =====================================================================
\section{Model}\label{sec:model}
% =====================================================================

We work in the usual synodic frame of the circular restricted three-body
problem, non-dimensionalised so that the primary separation, the total mass and
the mean motion are unity. The larger primary sits at $x=-\mu$ and the smaller
at $x=1-\mu$, with $\mu=M_2/(M_1+M_2)$. Distances from the primaries are
\begin{equation}
  r_1 = \lVert\bm{r}-\bm{r}_1\rVert, \qquad
  r_2 = \lVert\bm{r}-\bm{r}_2\rVert ,
\end{equation}
and on the segment between them $r_1=x+\mu$ and $r_2=(1-\mu)-x$, so that
\begin{equation}
  r_1 + r_2 = 1 \quad\text{identically.}
  \label{eq:unitsum}
\end{equation}
Equation~\eqref{eq:unitsum} is exact, not asymptotic, and is used in
Section~\ref{sec:closed}.

\subsection{Sail acceleration}

We model an ideal specular reflector. With $\alpha$ the cone angle between the
Sun line and the sail normal and $\delta$ the clock angle about the Sun line,
\begin{equation}
  \bm{a}_{\mathrm{sail}}
  = \beta\,\frac{1-\mu}{r_1^{2}}\,\cos^{2}\!\alpha\;\nhat ,
  \qquad
  \nhat = \cos\alpha\,\rhat + \sin\alpha\cos\delta\,\phat
        + \sin\alpha\sin\delta\,\qhat ,
  \label{eq:sail}
\end{equation}
where $\beta$, the lightness number, is the ratio of the peak radiation-pressure
acceleration to solar gravity. The force is directed along the normal with
magnitude proportional to $\cos^{2}\alpha$; this is the specular reflector, and
it is distinct from the perfect absorber, for which the force lies along the Sun
line with magnitude proportional to $\cos\alpha$. The two coincide only at
$\alpha=0$.

The triad in~\eqref{eq:sail} must be orthonormal. We take
\begin{equation}
  \rhat = \frac{\bm{r}-\bm{r}_1}{r_1},\qquad
  \phat = \frac{\hat{\bm{k}}\times\rhat}{\lVert\hat{\bm{k}}\times\rhat\rVert},
  \qquad
  \qhat = \rhat\times\phat ,
  \label{eq:triad}
\end{equation}
with $\hat{\bm{k}}$ the out-of-ecliptic unit vector. Substituting
$\hat{\bm{k}}$ itself for $\qhat$ --- a natural but incorrect shortcut --- gives
a non-orthonormal triad, since $\rhat\cdot\hat{\bm{k}}=z/r_1\neq0$ off the
ecliptic, and then
$\lVert\nhat\rVert^{2}=1+\sin(2\alpha)\sin\delta\,(z/r_1)$, so that the nominal
unit normal is not a unit vector and both the magnitude and the direction of the
force are wrong. The error is \SI{1.0}{\percent} at $z/r_1=0.02$ and
\SI{4.4}{\percent} at $z/r_1=0.09$. It vanishes identically at $\alpha=0$, so no
face-on result in this paper is affected; we record it because off-axis
quantities computed on the shortcut triad are not reliable.
Figure~\ref{fig:schematic} shows the geometry.

\begin{figure}[t]
  \centering
  \includegraphics[width=\textwidth]{fig1_schematic.png}
  \caption{Geometry of the face-on solar-sail CR3BP.
  \textbf{(a)} The synodic frame. The classical $L_1$ lies essentially on the
  Hill sphere ($r_2 = 0.9967\,r_H$); as $\beta$ grows the equilibrium migrates
  sunward, reaching the tidal-parity radius $r_2=\mu^{1/3}=3^{1/3}r_H$ at
  $\beta=0.0286$. The Earth neighbourhood is exaggerated by a factor of 18;
  at true scale every feature would collapse onto the smaller primary.
  \textbf{(b)} Attitude convention and the orthonormal
  triad~\eqref{eq:triad}. At $\alpha=0$ the normal is $\rhat$, the force is
  radial, and the sail is conservative.}
  \label{fig:schematic}
\end{figure}

\subsection{The face-on case is conservative}

Setting $\alpha=0$ in~\eqref{eq:sail} gives
$\bm{a}_{\mathrm{sail}}=\beta(1-\mu)r_1^{-2}\rhat$, a central force derivable
from the potential $-\beta(1-\mu)/r_1$. It therefore merely rescales the larger
primary, and the entire problem reduces to a CR3BP with modified pseudo-potential
\begin{equation}
  U = \tfrac{1}{2}\!\left(x^{2}+y^{2}\right)
      + \frac{(1-\beta)(1-\mu)}{r_1} + \frac{\mu}{r_2} .
  \label{eq:U}
\end{equation}
A Jacobi integral survives,
\begin{equation}
  C_{\mathrm{sail}} = 2U - v^{2},
  \label{eq:jacobi}
\end{equation}
and it is $U$ of~\eqref{eq:U}, not the unscaled gravitational pseudo-potential,
that must be used: the two differ by $2\beta(1-\mu)/r_1$, which varies along an
orbit, so the gravitational Jacobi constant is not an integral for $\beta\neq0$.
On a representative family member ($\beta=0.001$) we measure a spread of
\num{1.8e-15} in $C_{\mathrm{sail}}$ against \num{5.8e-6} in its gravitational
counterpart over one full period.

For $\alpha\neq0$ the force acquires a transverse component, is no longer
conservative, and no Jacobi integral exists. Every result in
Sections~\ref{sec:dissolution}--\ref{sec:families} is therefore restricted to
$\alpha=0$; Section~\ref{sec:control} is the sole exception and is explicitly a
statement about the linearisation in the attitude angles.

% =====================================================================
\section{Dissolution of the collinear structure}\label{sec:dissolution}
% =====================================================================

\subsection{The hovering law}

Deleting the smaller primary from~\eqref{eq:U} leaves an on-axis balance whose
root is
\begin{equation}
  x_{\mathrm{hover}}(\beta) = \bigl[(1-\beta)(1-\mu)\bigr]^{1/3},
  \label{eq:hover}
\end{equation}
the radius at which reduced solar gravity balances centrifugal acceleration at
the synchronous rate. This is a heliocentric hovering point, not a perturbed
Lagrange point. The full three-body equilibrium satisfies
\begin{equation}
  x = \frac{(1-\beta)(1-\mu)}{r_1^{2}} - \frac{\mu}{r_2^{2}} ,
  \label{eq:balance}
\end{equation}
and the difference between its root and~\eqref{eq:hover} is the smaller
primary's entire contribution. At $\beta=0.5$ that contribution is
\SI{3817}{\kilo\metre} against a standoff of \SI{30865449}{\kilo\metre}, i.e.
$20.6$ Hill radii: the Earth is dynamically irrelevant there, and results quoted
at such $\beta$ describe hovering geometry rather than three-body dynamics.

\subsection{Linear structure and the saddle strength}

Linearising~\eqref{eq:U} on the axis gives
$U_{xx}=1+2A$, $U_{yy}=1-A$, $U_{zz}=-A$ with
\begin{equation}
  A(\beta) = \underbrace{\frac{(1-\beta)(1-\mu)}{r_1^{3}}}_{\textstyle S}
           + \underbrace{\frac{\mu}{r_2^{3}}}_{\textstyle s} .
  \label{eq:A}
\end{equation}
This $A$ is not a new quantity. Setting $\beta=0$ and using
$r_1+r_2=1$, \eqref{eq:A} reduces to
$\mu/\gamma^{3}+(1-\mu)/(1-\gamma)^{3}$ with $\gamma=r_2$, which is exactly the
second Legendre coefficient $c_2$ of the collinear expansion introduced by
Richardson~\cite{richardson1980} and used throughout halo-orbit theory
\cite{gomez2001,ceccaroni2016}. We verify the identity numerically to machine
zero at Sun--Earth, Earth--Moon and Sun--Jupiter. We therefore write
\begin{equation}
  A \equiv c_2 ,
  \label{eq:Ac2}
\end{equation}
and use $c_2$ from here on. The sail's effect is to replace $(1-\mu)$ by
$(1-\beta)(1-\mu)$ in that coefficient: \emph{a solar sail continuously detunes
$c_2$}, which is otherwise fixed once the system and the libration point are
chosen. This is the compact statement of what a sail contributes, and we return
to it in Section~\ref{sec:freq}.

The out-of-plane motion is a pure oscillator, $\nu=\sqrt{c_2}$. The in-plane
characteristic polynomial is
\begin{equation}
  \lambda^{4} + (2-A)\lambda^{2} + \bigl(1+A-2A^{2}\bigr) = 0 ,
  \label{eq:quartic}
\end{equation}
whose roots are $\lambda^{2}=\tfrac12\bigl[(A-2)\pm D\bigr]$ with
$D=\sqrt{A(9A-8)}$, giving a saddle exponent $\lambda_u$ and a centre frequency
$\omega$,
\begin{equation}
  \lambda_u^{2} = \tfrac12\bigl[(A-2)+D\bigr], \qquad
  \omega^{2} = \tfrac12\bigl[D-(A-2)\bigr].
  \label{eq:rates}
\end{equation}
A real hyperbolic pair requires $c_2>1$.

As the equilibrium recedes from the smaller primary, $r_1^3\to(1-\beta)(1-\mu)$
by~\eqref{eq:hover} and hence $S\to1$, so
\begin{equation}
  c_2 - 1 \;\longrightarrow\; s = \frac{\mu}{r_2^{3}} .
  \label{eq:As}
\end{equation}
We call $s$ the \emph{saddle strength}. Because $s>0$ for all finite $r_2$, the
saddle never strictly disappears: there is no bifurcation, only a smooth decay,
and the physically meaningful statement is its rate. Table~\ref{tab:collapse}
records that rate; $s$ falls four decades between $\beta=0$ and $\beta=0.5$.

It is worth being precise about~\eqref{eq:As}, because the convergence is slower
than it looks. At tidal parity, $S=(1-\mu)/(1-\mu^{1/3})=1.014635$ exactly, so
$c_2=2.014635$ rather than $2$, and $\lambda_u=1.505418$ rather than
$5^{1/4}=1.495349$. The identification $c_2-1=s$ carries an $O(\mu^{1/3})$ error,
\SI{0.73}{\percent} in $c_2$ at the Sun--Earth threshold. Both terms
in~\eqref{eq:A} are nonetheless of smaller-primary origin: with the primary
deleted, $S=1$ identically, so $S-1$ measures its indirect effect through the
displacement of the equilibrium, while $s$ measures the direct tidal term. We
therefore write $A-1=(S-1)+s$ and use the exact $A$ throughout.

\begin{table}[t]
  \centering
  \caption{Collapse of the saddle strength $s=\mu/r_2^{3}$ and of the unstable
  eigenvalue with sail lightness number, Sun--Earth
  ($\mu=\num{3.003e-6}$). Tidal parity is $s=1$.}
  \label{tab:collapse}
  \begin{tabular}{lccc}
    \toprule
    $\beta$ & $s$ & $\lambda_u$ [nd] & note \\
    \midrule
    $0$        & \num{3.0303}   & \num{2.53256} & classical $L_1$ \\
    $0.010$    & \num{2.123}    & \num{2.13831} & \\
    $0.028646$ & \num{1.000}    & \num{1.50542} & tidal parity \\
    $0.050$    & \num{0.4012}   & \num{0.99449} & \\
    $0.100$    & \num{0.06820}  & \num{0.44400} & \\
    $0.500$    & \num{3.419e-4} & \num{0.03578} & $20.6\,r_H$ standoff \\
    \bottomrule
  \end{tabular}
\end{table}

\begin{figure}[t]
  \centering
  \includegraphics[width=\textwidth]{fig8_structure_dissolution.png}
  \caption{Dissolution of the Sun--Earth collinear structure.
  \textbf{(a)} The full three-body equilibrium is indistinguishable from the
  Earth-free hovering law~\eqref{eq:hover}; the inset shows the smaller
  primary's entire contribution collapsing to a few thousand kilometres.
  \textbf{(b)} The standoff crosses the Hill sphere almost immediately
  ($\beta_{\mathrm{crit}}=\num{2.98e-4}$) and reaches $20.6\,r_H$ by
  $\beta=0.5$. \textbf{(c)} The saddle strength falls four decades; $s=1$
  locates tidal parity. \textbf{(d)} Both frequencies converge on the mean
  motion and $\lambda_u\to0$: the Keplerian epicyclic degeneracy of an inverse
  square field, not a resonance with the smaller primary.}
  \label{fig:dissolution}
\end{figure}

Two consequences of Figure~\ref{fig:dissolution}(d) deserve emphasis, because
both have been misread in the past. As $c_2\to1$ the planar roots
of~\eqref{eq:quartic} become $\lambda^{2}\in\{0,-1\}$ and $\nu=\sqrt{A}\to1$,
so the in-plane period is driven to exactly $2\pi$ in non-dimensional units ---
one year for Sun--Earth. This is \emph{forced}, not resonant: in an inverse
square field the epicyclic frequency equals the mean motion identically.
Second, $\lambda_u\to0$ does not indicate that the sail has stabilised a saddle;
the orbit has simply moved away from one, and the residual instability is
precisely the smaller primary's leftover tidal term.

% =====================================================================
\section{The tidal-parity threshold in closed form}\label{sec:closed}
% =====================================================================

\subsection{Derivation}

Define tidal parity as the condition that the direct tidal term equal the
centrifugal background of the synodic frame,
\begin{equation}
  s = \frac{\mu}{r_2^{3}} = 1
  \qquad\Longleftrightarrow\qquad
  r_2 = \mu^{1/3} = 3^{1/3}\,r_H ,
  \label{eq:parity}
\end{equation}
where $r_H=(\mu/3)^{1/3}$. Below this radius the smaller primary dominates the
Hessian; above it the centrifugal and rescaled-solar terms do.

Impose~\eqref{eq:parity} on the balance~\eqref{eq:balance}. By the exact
identity~\eqref{eq:unitsum}, $r_2=\mu^{1/3}$ forces
\begin{equation}
  r_1 = 1-\mu^{1/3}
\end{equation}
with no error term. Substituting, and using $x=(1-\mu)-r_2$ and
$\mu/r_2^{2}=\mu^{1/3}$,
\begin{equation}
  \underbrace{(1-\mu)-\mu^{1/3}}_{x}
  = \frac{(1-\beta)(1-\mu)}{\bigl(1-\mu^{1/3}\bigr)^{2}} - \mu^{1/3} .
\end{equation}
The $\mu^{1/3}$ terms cancel identically, leaving
$(1-\mu)=(1-\beta)(1-\mu)\bigl(1-\mu^{1/3}\bigr)^{-2}$ and hence
$1-\beta=\bigl(1-\mu^{1/3}\bigr)^{2}$. Therefore
\begin{equation}
  \boxed{\;
  \beta_{\mathrm{crit}} = 1-\bigl(1-\mu^{1/3}\bigr)^{2}
                        = 2\mu^{1/3}-\mu^{2/3}
  \;}
  \label{eq:closed}
\end{equation}
This is exact for every $\mu$; the cancellation is algebraic and involves no
expansion in $\mu$. To leading order $\beta_{\mathrm{crit}}\simeq2\mu^{1/3}$.

The same substitution gives the linear structure at parity in closed form,
\begin{equation}
  c_2^{\mathrm{parity}} = 1 + \frac{1-\mu}{1-\mu^{1/3}} ,
  \label{eq:Aparity}
\end{equation}
which for Sun--Earth is $2.014635$, the value quoted in
Section~\ref{sec:dissolution}.

\subsection{Verification}

We evaluate~\eqref{eq:closed} against independent Brent root-finding on
$s(\beta)-1$, where $s$ is computed from the numerically solved
equilibrium~\eqref{eq:balance}. Table~\ref{tab:closed} reports both across five
decades of mass ratio. The worst discrepancy over a 61-point logarithmic sweep
of $\mu\in[10^{-7},10^{-2}]$ is \num{2.1e-16}, i.e. machine precision; the two
computations share no code path beyond the balance itself.

\begin{table}[t]
  \centering
  \caption{The closed form~\eqref{eq:closed} against independent root-finding on
  $s(\beta)=1$. Agreement is at machine precision for every mass ratio tested.
  $r_2/r_H=3^{1/3}=1.4422$ at parity by construction, independent of $\mu$.}
  \label{tab:closed}
  \begin{tabular}{lccccc}
    \toprule
    system & $\mu$ &
    $\beta_{\mathrm{crit}}$ (closed form) &
    $\beta_{\mathrm{crit}}$ (root-find) &
    $|\Delta|$ & $r_2$ [nd] \\
    \midrule
    Sun--Mercury & \num{1.660e-7} & \num{0.010961524792} & \num{0.010961524792} & \num{1.7e-16} & \num{0.005500} \\
    Sun--Earth   & \num{3.003e-6} & \num{0.028646456169} & \num{0.028646456169} & \num{6.9e-18} & \num{0.014427} \\
    Sun--Jupiter & \num{9.537e-4} & \num{0.187175530204} & \num{0.187175530204} & \num{5.6e-17} & \num{0.098438} \\
    Earth--Moon  & \num{1.215e-2} & \num{0.406934946280} & \num{0.406934946280} & \num{5.6e-17} & \num{0.229900} \\
    \bottomrule
  \end{tabular}
\end{table}

For Sun--Earth the threshold corresponds to a standoff
$r_2=\num{0.014427}\,\nd=\SI{2158294}{\kilo\metre}=1.4422\,r_H$. For reference,
the classical $L_1$ already sits at $r_2=0.9967\,r_H$, essentially on the Hill
sphere, so the Hill-sphere exit criterion $r_2=r_H$ is reached almost
immediately, at $\beta=\num{2.9814e-4}$, and is not a useful discriminant.

% =====================================================================
\section{The frequency ratio: a universal bound on \texorpdfstring{$c_2$}{c2}}\label{sec:freq}
% =====================================================================

The threshold~\eqref{eq:closed} is a level crossing of a smooth monotone
function. Nothing bifurcates there, and the choice of the level $s=1$ is partly
a statement about the non-dimensionalisation. The standard route to promoting
such a marker into a dynamical event is a commensurability of the linear
frequencies. This section searches for one and reports that none exists ---
together with an exact extremum that does.

\subsection{The bound and its extremum}

From~\eqref{eq:rates}, the ratio of the out-of-plane to in-plane frequency
depends on $c_2$ alone,
\begin{equation}
  \frac{\nu^{2}}{\omega^{2}} = \frac{2c_2}{D-(c_2-2)},
  \qquad D=\sqrt{c_2(9c_2-8)} .
  \label{eq:R}
\end{equation}
The two frequencies in~\eqref{eq:rates} are standard --- they appear in this
form in~\cite{gomez2001} and~\cite{ceccaroni2016} --- but their ratio is, to our
knowledge, not formed anywhere. Doing so is what follows.

Differentiating~\eqref{eq:R}, the stationarity condition reduces to $D=2c_2$,
i.e. $9c_2^{2}-8c_2=4c_2^{2}$, giving
\begin{equation}
  \boxed{\;c_2^\star = \tfrac{8}{5}\;}
  \label{eq:Astar}
\end{equation}
and there, from~\eqref{eq:rates},
\begin{equation}
  \lambda_u = \sqrt{\tfrac{7}{5}},\qquad
  \nu = \sqrt{\tfrac{8}{5}},\qquad
  \omega = \sqrt{\tfrac{9}{5}}
  \qquad\Longrightarrow\qquad
  \lambda_u^{2}:\nu^{2}:\omega^{2} = 7:8:9 ,
  \label{eq:789}
\end{equation}
\begin{equation}
  \left.\frac{\nu}{\omega}\right|_{c_2^\star} = \sqrt{\tfrac{8}{9}}
  = \frac{2\sqrt2}{3} = 0.942809041582\ldots
  \label{eq:bound}
\end{equation}
Over the whole physically admissible range $c_2>1$ the ratio is confined to
\begin{equation}
  \frac{\nu}{\omega} \in \Bigl[\tfrac{2\sqrt2}{3},\,1\Bigr) ,
  \label{eq:band}
\end{equation}
a band \SI{5.72}{\percent} wide, approaching unity in the hovering limit
$c_2\to1$ and again as $c_2\to\infty$. Numerically, a $4\times10^{5}$-point scan
places the minimum at $c_2=1.60000000$ with value $0.942809041582$, agreeing
with~\eqref{eq:Astar} and~\eqref{eq:bound} to \num{8.1e-10} and \num{1.1e-16}
respectively. Note that $c_2^\star$, the ratios~\eqref{eq:789} and the
bound~\eqref{eq:bound} contain no reference to $\mu$: they are properties of the
collinear linearisation itself, shared by every three-body system.

The narrowness of~\eqref{eq:band} has been observed qualitatively. Ceccaroni
\emph{et al.}~\cite{ceccaroni2016}, sweeping all three collinear points over the
full range $\mu\in(0,\tfrac12]$, note that the centre-manifold dynamics is
``characterized by almost equal values of the frequencies for all mass
ratios''. Equations~\eqref{eq:band} and~\eqref{eq:bound} make that observation
exact and supply the constant it is missing.

\subsection{No low-order resonance is attainable}

We tested fourteen low-order rationals --- $1{:}2$, $2{:}3$, $3{:}4$, $4{:}5$,
$5{:}6$, $8{:}9$, $1{:}1$, $9{:}8$, $6{:}5$, $5{:}4$, $4{:}3$, $3{:}2$,
$5{:}3$, $2{:}1$ --- against the band~\eqref{eq:band}. \emph{None} lies inside
it. Since $\nu<\omega$ strictly for $c_2>1$, even the $1{:}1$ commensurability is
approached only asymptotically, as the equilibrium degenerates into the hovering
point.

The consequence is that no commensurability underwrites tidal parity, and none
exists elsewhere in the face-on solar-sail CR3BP either. Equation
\eqref{eq:closed} must therefore be presented as an exact \emph{scale marker}
with a clean geometric reading, not as a dynamical transition. We regard this
negative as a result rather than an omission: it is the check a reader is
entitled to demand, and it forecloses an over-claim.

Two caveats. First, \eqref{eq:band} constrains the \emph{linear} frequencies
only. Halo families bifurcate from planar Lyapunov families at a
commensurability of the \emph{nonlinear} frequencies, reached at finite
amplitude; the classical $L_1$ has $\nu/\omega=0.9659$ and halo orbits exist
there regardless. Nothing here constrains halo existence.
Second, $c_2^\star$ is a stationary point, not a transition: the dynamics change
character nowhere in $c_2$. Its value is that a stationary point is intrinsic ---
it survives any smooth reparametrisation and any choice of units --- whereas a
level crossing at $s=1$ does not.

\subsection{What the sail contributes}

It would be natural to suppose $c_2^\star=8/5$ inaccessible without continuous
thrust, on the grounds that the classical collinear points have $c_2\geq4$. That
is true only of $L_1$. Computing $c_2$ at all three classical collinear points
over $\mu\in(0,\tfrac12]$ gives
\begin{equation}
  c_2(L_1)\in[4.042,\,8.000], \qquad
  c_2(L_2)\in[1.570,\,3.959], \qquad
  c_2(L_3)\in[1.000,\,1.570],
\end{equation}
so the classical problem covers $c_2<4$ freely at $L_2$ and $L_3$, and it
\emph{does} attain $c_2^\star$: at $L_2$ with $\mu=0.480187660$, where
$\nu/\omega$ equals~\eqref{eq:bound} to twelve digits. $L_3$ at $\mu=\tfrac12$
reaches $c_2=1.569787$ and misses the bound by only \num{3.07e-5} --- close
enough that the near-coincidence is worth recording, since $L_3$ occupies
precisely the interval in which the extremum lies.

The sail's contribution is therefore \emph{tunability}, not exclusivity, and
\eqref{eq:Ac2} is the compact way to say it: at fixed $\mu$ a sail sweeps $c_2$
continuously through $(1,c_2^{\mathrm{classical}}]$, whereas classically $c_2$
is fixed once the system and the point are chosen. Sun--Earth classically offers
only $c_2=4.060819$ at $L_1$, $3.940764$ at $L_2$ and $1.000003$ at $L_3$; none
is at the bound, and a sail at $\beta=0.040932$ places that system exactly on it
--- a factor $1.4289$ beyond tidal parity, at $s=0.5897$.
Figure~\ref{fig:freq} collects all of this.

\begin{figure}[t]
  \centering
  \includegraphics[width=\textwidth]{fig4_frequency_ratio.png}
  \caption{The frequency ratio at collinear equilibria.
  \textbf{(a)} $\nu/\omega$ against $c_2$, with the exact minimum
  $2\sqrt2/3$ at $c_2=8/5$ and, in the reserved strip below, the ranges of $c_2$
  reached by the three classical collinear points over $\mu\in(0,\tfrac12]$;
  the vertical placement of those bars is a layout choice and carries no value.
  \textbf{(b)} The three linear rates, meeting the ratio $7:8:9$ at
  $c_2^\star$. \textbf{(c)} The sail's tuning curve for Sun--Earth:
  $c_2(\beta)$ sweeps continuously from the classical value down to unity,
  passing tidal parity and then $c_2^\star$. Open markers give the three
  classical Sun--Earth values, none of which is at the bound.
  \textbf{(d)} Classical $c_2$ against $\mu$ for all three points;
  $c_2^\star$ is attained at $L_2$ for $\mu=0.4802$.}
  \label{fig:freq}
\end{figure}

% =====================================================================
\section{Halo families across the flown sail band}\label{sec:families}
% =====================================================================

The band $\beta\in[0.001,0.05]$ is where a sail both matters dynamically and
still inhabits a genuine three-body structure: below it the sail is irrelevant,
and above tidal parity the collinear structure has lost its meaning. We
therefore walk the halo family at twelve values of $\beta$ across that band.

Two methodological points are load-bearing. First, all lengths are expressed in
units of the local scale $\gamma=\lvert(1-\mu)-x_{\mathrm{eq}}\rvert$, which
doubles across the band ($0.0101$ at $\beta=0.001$ to $0.0196$ at $\beta=0.05$).
Absolute step sizes and amplitude bounds therefore mean different things at
different $\beta$; fixed ones admit far-field orbits, and a seed amplitude of
$A_z=0.02$ --- reasonable for Earth--Moon --- is roughly $2\gamma$ for
Sun--Earth, larger than the entire distance from the equilibrium to the smaller
primary.

Second, the differential corrector solves for $[x_0,\,\dot y_0,\,T/2]$ subject
to $[\dot x_f,\,\dot z_f,\,y_f]=0$; $z_0$ is not a free variable and $A_z$ is
not a constraint. Stepping $A_z$ and re-solving therefore fails at folds in
$A_z$, where the branch turns and no solution exists at the requested
amplitude. We use pseudo-arclength continuation, appending the arclength
condition and treating the continuation parameter as an unknown, so that folds
are traversed rather than tripped over. Branch identity is additionally
enforced across $\beta$: each $\beta$ is seeded from its predecessor's converged
state with $z_0$ rescaled by $\gamma_{\mathrm{new}}/\gamma_{\mathrm{old}}$, and
two guards verify the result --- the sign of the $z$-extremum asymmetry
\begin{equation}
  \Delta_z = \frac{\lvert z_0\rvert - \lvert z(T/2)\rvert}
                  {\lvert z_0\rvert + \lvert z(T/2)\rvert}
\end{equation}
must match the reference $\beta$, and the family must lie sunward of the
equilibrium. $\Delta_z$ is $O(0.1)$ for a genuine halo and vanishes identically
on the vertical-Lyapunov branch, which makes it an effective discriminator.
Without the cross-$\beta$ guard one of the twelve families lands on a different
branch.

Figure~\ref{fig:atlas} reports the resulting atlas. The Jacobi constant is
computed from~\eqref{eq:jacobi}; because the potential~\eqref{eq:U} is itself
$\beta$-dependent, raw $C_{\mathrm{sail}}$ is not comparable across $\beta$, and
we plot instead the depth below the local equilibrium,
$C_{\mathrm{sail}}-C_{\mathrm{sail}}(x_{\mathrm{eq}})$, against $A_z/\gamma$.

\begin{figure}[t]
  \centering
  \includegraphics[width=\textwidth]{fig10_halo_atlas.png}
  \caption{Halo families across $\beta\in[0.001,0.05]$, Sun--Earth,
  $\alpha=0$, walked by pseudo-arclength continuation with branch identity
  enforced across $\beta$. \textbf{(a)} Family depth below the local
  equilibrium in the sail potential, against amplitude scaled by the local
  length $\gamma$. \textbf{(b)} Period. \textbf{(c)} Stability index, with the
  $\lvert\nu\rvert=1$ boundary. \textbf{(d)} The $(\beta,A_z)$ plane: family
  extent, folds, and linearly stable segments, with tidal parity marked.}
  \label{fig:atlas}
\end{figure}

\begin{figure}[t]
  \centering
  \includegraphics[width=0.78\textwidth]{fig2_stability.png}
  \caption{Collapse of the unstable eigenvalue with $\beta$. The trend is not
  sail stabilisation of a saddle: the standoff grows, the smaller primary's
  tidal term dies, and the orbit moves away from the saddle it began near.}
  \label{fig:stability}
\end{figure}

% =====================================================================
\section{Control authority of a face-on sail}\label{sec:control}
% =====================================================================

The dynamics above are open-loop. Since equilibria in this region are unstable,
it is natural to ask what a sail can do about it, and here the face-on
configuration turns out to be singular.

A sail has two control inputs, the attitude angles, and its thrust direction is
rigidly coupled to its magnitude through~\eqref{eq:sail}. Linearising in
$(\alpha,\delta)$ about a nominal $(\alpha_0,\delta_0)$ gives a $6\times2$
control Jacobian $B_{\mathrm{sail}}$ whose non-trivial rows are
$\partial\bm{a}/\partial\alpha$ and $\partial\bm{a}/\partial\delta$. At
$\alpha_0=0$ exactly,
\begin{equation}
  \frac{\partial\bm{a}}{\partial\alpha}
  = \beta\frac{1-\mu}{r_1^{2}}
    \bigl(\cos\delta_0\,\phat + \sin\delta_0\,\qhat\bigr),
  \qquad
  \frac{\partial\bm{a}}{\partial\delta} = \bm{0} .
  \label{eq:Bsail}
\end{equation}
The clock angle has \emph{no} first-order authority when the sail is face-on: it
only rotates a normal already parallel to $\rhat$. The single available
direction is therefore fixed by the constant $\delta_0$, and it is purely
transverse, since
$\bm{a}\cdot\rhat=\beta(1-\mu)r_1^{-2}\cos^{3}\alpha$ and
$\mathrm{d}(\cos^{3}\alpha)/\mathrm{d}\alpha\to0$ as $\alpha\to0$.

Rank of $B$ is not the relevant test, because Coriolis coupling can carry a
transverse input into the radial direction. Applying the Kalman rank condition
to the linearised system at the $\beta=0.05$ equilibrium ($c_2=1.409194$):

\begin{center}
\begin{tabular}{lcc}
  \toprule
  control model & Kalman rank & verdict \\
  \midrule
  unconstrained three-axis thruster ($6\times3$) & $6/6$ & controllable \\
  sail, $\alpha_0=0$, $\delta_0=0$   & $4/6$ & \textbf{uncontrollable} \\
  sail, $\alpha_0=0$, $\delta_0=\pi/2$ & $2/6$ & \textbf{uncontrollable} \\
  sail, $\alpha_0\geq0.5^\circ$      & $6/6$ & controllable \\
  \bottomrule
\end{tabular}
\end{center}

The in-plane four-state block \emph{is} controllable by a single transverse
input (rank $4/4$) --- Coriolis does the work --- so the absent radial authority
is not by itself fatal. What fails is the out-of-plane mode: $z$ decouples at
linear order, and because $\partial\bm{a}/\partial\delta=\bm{0}$ the sail cannot
redirect thrust into it without changing $\alpha_0$. A face-on sail reaches the
in-plane pair or the vertical mode, never both:
\begin{equation}
  \alpha_0 = 0 \text{ is a singular nominal for sail station-keeping.}
\end{equation}
Any non-zero cone angle restores full rank; $0.5^\circ$ suffices formally,
though the second singular value there is \num{4.5e-4} against \num{5.2e-2} for
the first, a $115{:}1$ conditioning penalty, so usable authority requires a few
degrees.

Authority is maximised at
\begin{equation}
  \alpha_0^\star = \arctan\!\bigl(1/\sqrt2\bigr) = 35.264^\circ,
  \qquad
  \frac{\sigma_2}{\sigma_1} = \frac13,
  \qquad
  \cos^{2}\alpha_0^\star = \frac23 ,
  \label{eq:alphastar}
\end{equation}
all exact and independent of $\beta$. This angle is the classical optimal sail
cone angle for maximum transverse thrust~\cite[\S2.6]{mcinnes1999}; recovering
it from the control Jacobian is a benchmark check on~\eqref{eq:Bsail}, not a new
result. What is new is that the same angle maximises out-of-plane control
authority, with $\sigma_2/\sigma_1$ exactly $1/3$ there, and that the authority
vanishes identically at $\alpha_0=0$. Figure~\ref{fig:authority} shows the
trade: thrust falls as $\cos^{2}\alpha_0$ while authority rises from zero, and
$\alpha_0^\star$ is where the second is maximal.

\begin{figure}[t]
  \centering
  \includegraphics[width=\textwidth]{fig5_control_authority.png}
  \caption{Control authority of a solar sail at the artificial equilibrium,
  $\beta=0.05$. \textbf{(a)} Singular values of $B_{\mathrm{sail}}$ against
  nominal cone angle; $\sigma_2\to0$ as $\alpha_0\to0$. \textbf{(b)}
  Conditioning, maximal at $\alpha_0^\star$ with the exact value $1/3$; the
  shaded strip at the origin is the uncontrollable nominal. \textbf{(c)}
  Out-of-plane authority scales linearly with $\beta$. \textbf{(d)} The design
  trade, both curves dimensionless: thrust $\cos^{2}\alpha_0$ against
  normalised authority.}
  \label{fig:authority}
\end{figure}

% =====================================================================
\section{Where the threshold sits against flown hardware}\label{sec:tech}
% =====================================================================

The threshold~\eqref{eq:closed} is only interesting if it is within reach. We
therefore reduce $\beta$ from primary mission specifications rather than quoting
secondary values, using $\beta=\sigma^{\ast}/\sigma$ with the critical sail
loading $\sigma^{\ast}=2P/g_{\odot}$. Evaluated from CODATA and IAU constants,
$\sigma^{\ast}=\SI{1.5311}{\gram\per\metre\squared}$, reproducing McInnes'
\num{1.53}~\cite{mcinnes1999} to four figures, which validates the chain.

\begin{table}[t]
  \centering
  \caption{Lightness number of every solar sail flown or funded to flight
  hardware, reduced from primary specifications. $\beta_{\mathrm{ideal}}$ is the
  perfect-reflector limit; $\beta_{\mathrm{eff}}$ applies the optical efficiency
  $\eta=1.12/1.79=0.626$ measured on IKAROS. IKAROS is the only entry whose
  lightness number follows from a \emph{measured} thrust and needs no optical
  assumption.}
  \label{tab:tech}
  \begin{tabular}{lccccl}
    \toprule
    mission & area [\si{\metre\squared}] & mass [\si{\kilogram}] &
    $\beta_{\mathrm{ideal}}$ & $\beta_{\mathrm{eff}}$ & status \\
    \midrule
    IKAROS~\cite{mori2010,tsuda2011}   & 196  & 307 & \num{0.00098} & \num{0.00062} & flown 2010, measured \\
    LightSail-2~\cite{spencer2021}     & 32   & 5   & \num{0.00980} & \num{0.00613} & flown 2019 \\
    NEA Scout~\cite{mcnutt2014}        & 86   & 14  & \num{0.00941} & \num{0.00588} & sail not deployed, lost 2022 \\
    ACS3~\cite{acs3}                   & 80   & 16  & \num{0.00766} & \num{0.00479} & flown 2024 \\
    Solar Cruiser                      & 1653 & --- & \num{0.02024} & ---           & cancelled 2022 \\
    \bottomrule
  \end{tabular}
\end{table}

No sail flown and deployed exceeds $\beta=\num{0.00613}$, and the only measured
value is IKAROS at $\beta=\num{0.00062}$. Against
$\beta_{\mathrm{crit}}=\num{0.028646}$:

\begin{center}
\begin{tabular}{lcc}
  \toprule
  reference & $\beta$ & shortfall \\
  \midrule
  best flown, perfect-reflector limit (LightSail-2) & \num{0.009799} & $2.9\times$ \\
  best flown, realistic (LightSail-2)               & \num{0.006131} & $4.7\times$ \\
  only measured value (IKAROS)                      & \num{0.000615} & $47\times$ \\
  most ambitious funded design (Solar Cruiser, cancelled) & \num{0.020236} & $1.4\times$ \\
  \bottomrule
\end{tabular}
\end{center}

Tidal parity is therefore roughly three times beyond the best sail flown and
within \SI{40}{\percent} of the most ambitious design funded to flight hardware.
The defensible statement is that it lies at the edge of \emph{near-term} rather
than current capability. We note explicitly that a claim of the form
``$\beta\sim0.01$--$0.05$ is current technology'', which appears informally, is
not supported by any flown mission.

A consequence worth stating: at the best flown $\beta\approx0.006$ the saddle
strength is $s\approx2.5$, so for every sail ever flown the collinear structure
is fully intact and smaller-primary-dominated. The dissolution characterised
here is a prediction about a regime not yet reached, not an account of observed
behaviour.

% =====================================================================
\section{Limitations}\label{sec:limits}
% =====================================================================

\begin{enumerate}
\item \textbf{Face-on only.} Everything except Section~\ref{sec:control} assumes
$\alpha=0$. For $\alpha\neq0$ the force is non-conservative, no Jacobi integral
exists, and the effective-potential analysis used throughout fails. The steered
case is open and is the natural sequel.

\item \textbf{Ideal specular reflector.} Membrane billow, finite slew rate,
self-shadowing, and wavelength-dependent optics are not modelled. Dachwald
\emph{et al.}~\cite{dachwald2013} argue persuasively that non-ideal optics matter
more than order-of-magnitude estimates suggest; the IKAROS efficiency
$\eta=0.626$ in Table~\ref{tab:tech} is a single-mission proxy, not a general
correction.

\item \textbf{No bifurcation is claimed.} $s>0$ for all finite $r_2$, so the
saddle persists; $c_2^\star$ is a stationary point of a smooth function. Both
thresholds are exact structural landmarks, and the resonance search of
Section~\ref{sec:freq} is what licenses that restraint.

\item \textbf{Derived hardware figures.} Every $\beta$ in Table~\ref{tab:tech}
except IKAROS' is reduced by us from published area and mass, not quoted by the
mission teams. Where a team publishes its own lightness number, prefer it.

\item \textbf{Linear controllability only.} Section~\ref{sec:control} is a
statement about the linearisation at the equilibrium. Nonlinear or
higher-order control strategies at $\alpha_0=0$ are not excluded; only
first-order authority is.
\end{enumerate}

% =====================================================================
\section{Conclusions}
% =====================================================================

A face-on solar sail rescales solar gravity rather than perturbing the collinear
point, so the collinear structure does not shift with $\beta$ --- it decays. We
have characterised that decay through the saddle strength $s=\mu/r_2^3$, which
falls four decades between $\beta=0$ and $\beta=0.5$, and obtained two exact
results.

The tidal-parity threshold admits the closed form
$\beta_{\mathrm{crit}}=2\mu^{1/3}-\mu^{2/3}$, exact for every mass ratio and
verified to machine precision across five decades of $\mu$. For Sun--Earth it is
$0.0286$, corresponding to a standoff of $3^{1/3}r_H$. The scale $\mu^{1/3}$ is
not itself new; the exact sail parameter required to reach it is.

The parameter $A$ of the collinear linearisation is exactly Richardson's
second Legendre coefficient $c_2$, and a sail's effect is to detune it
continuously. The frequency ratio at any collinear equilibrium is confined to
$\nu/\omega\in[2\sqrt2/3,1)$ and attains its bound at $c_2=8/5$, where the
linear rates stand in the ratio $7:8:9$. Both are independent of $\mu$. The
near-equality of the two frequencies has been noted
before~\cite{ceccaroni2016}; the bound and its extremum make it exact. No
low-order rational lies in the band, so no commensurability underwrites either
threshold, and we present them as scale markers rather than transitions. The
sail's role is to tune $c_2$ within a given system, reaching values ---
including the bound, at $\beta=0.0409$ for Sun--Earth --- that that system's
classical collinear points do not offer.

Finally, the face-on nominal is singular for station-keeping: the clock angle
has no first-order authority, the Kalman rank is $4/6$, and a few degrees of
cone angle are required to recover it. Authority peaks at the classical optimal
cone angle $\arctan(1/\sqrt2)$, where the conditioning is exactly $1/3$.

Set against flown hardware, tidal parity is a factor $4.7$ beyond the best
deployed sail and $47$ beyond the only measured lightness number. It is a
near-term design boundary, not a present one.

\section*{Data and code availability}

All numerical results, figures and the verification suite are reproducible from
a single entry point. Every quantity quoted in this paper is regenerated by the
pipeline and written to a machine-generated results file; no value is
hand-transcribed.

\bibliographystyle{unsrt}
\bibliography{refs}

\end{document}
